\section{Introduction}
\label{sec:introduction}

Internet connectivity is an input into an increasingly large set of production processes. For software developers, the dependence is especially direct: source code can be written locally, but modern development also relies on remote repositories, package registries, code review, issue tracking, continuous integration, documentation, and communication with collaborators. A nationwide loss of connectivity therefore interrupts both the transmission of code and many complementary tasks required to coordinate and ship software.

This paper studies a sharp episode of lost connectivity: Bangladesh's nationwide Internet shutdown in July 2024. Cloudflare reported that mobile networks began disconnecting early on July 18 and that national traffic fell to near zero later that day. Fixed broadband began to return on July 23, initially for selected sectors, while mobile connectivity recovered later, on July 28 \citep{belson2024}. Internet Society Pulse records the national disruption as beginning on July 18 and lasting roughly five and a half days \citep{isoc2024}. The event offers an unusual setting in which a digitally intensive form of work can be observed at high frequency before, during, and after a near-complete national loss of Internet connectivity.

The paper makes four contributions. First, an explicit international synthetic counterfactual measures the immediate collapse in observable public activity. Second, separate country-domain controls assess heterogeneity among users grouped by their most active pre-shutdown repository; the labels describe project applications, not verified employment industries. Third, Git metadata distinguish timestamped authorship from publication, revealing a small observed backlog alongside a large fall in the author-time creation of code eventually observed by August 10. A case review of residual authors and publishers documents selection on prior activity and a substantial automated component of continued publication. Fourth, cumulative national gaps and the separate within-Bangladesh recovery analysis describe persistence, without interpreting post-restoration baseline crossings as causal gains.

The empirical challenge is that GitHub activity is volatile and varies widely across countries and developers. I therefore construct a multi-country donor pool and use synthetic control as the primary design. The treated unit is Bangladesh, and the donor pool contains fourteen countries of different sample sizes, including 1,373 eligible Ghana-associated accounts and 446 Jordan-associated accounts; rates for the smaller donors are correspondingly noisier. The synthetic-control weights are chosen to match the pre-shutdown path of the daily active-user rate. The resulting counterfactual assigns 47.6 percent weight to Pakistan, 32.5 percent to Ghana, and 19.9 percent to Kenya. During July 18--23, the share of eligible Bangladeshi users with productive GitHub activity averaged only 2.02 percent, compared with 15.28 percent in the synthetic control. The estimated effect is $-$13.26 percentage points, or $-$86.8 percent relative to the synthetic counterfactual.

Several complementary designs support the same conclusion. An international difference-in-differences comparison yields an effect of $-$13.09 percentage points, and excluding national holidays changes the estimate only modestly to $-$13.38 percentage points. Holding the primary donor weights fixed, productive activity falls by 86.6 percent, shipping activity by 93.3 percent, and collaborative activity by 93.5 percent; re-optimizing the weights for each outcome gives similar declines of 86.7, 90.5, and 91.9 percent. Regression discontinuity in time is retained as an appendix sensitivity check.

The mechanism analysis separates observable online activity from underlying code creation. A Git commit can be created without Internet access, but pushing it to GitHub requires connectivity. For a large sample of public commits, I recover Git author and committer timestamps and match each commit to the first public PushEvent in which it appears. In the strict sample, in which the commit author is the same GitHub account that performs the push, 829 of the 2,932 commits authored during July 18--23 (28.3 percent) first appear in a public push on or after July 24; the committer timestamp gives 27.2 percent. Long publication delays are otherwise rare: 29.7 percent of shutdown-authored commits were pushed more than 24 hours after authorship, compared with an average of 1.1 percent across six pre-shutdown placebo windows of six days each. In levels, however, the backlog is small. For the cohort authored after the nationwide cutoff at approximately 15:00 UTC on July 18, observable author-time commit creation is 96--97 percent below comparable pre-shutdown windows (96.35 percent against duration-matched windows), and that cohort's delayed commits amount to 2.50 percent of normal volume (\cref{sec:backlog}). These counts cover only code pushed publicly by August 10. Later pushes, unresolved metadata, and code never published are unobserved, so the data cannot separate work that was not done from work published after August 10 or never published.

Residual publication is especially informative about measurement. Only 40 eligible accounts record a primary event on July 19--22. Individual examination identifies automated heads in 113 of their 180 pushes, including dependency updates, profile graphics, metrics and deployment output. Removing those pushes alone reduces daily residual activity from 0.212 to 0.117 percent, while preserving the original national specification. These accounts overlap only partly with the 250 strict onset authors; neither continued posting nor matching Git author metadata establishes that a person kept coding inside Bangladesh.

A separate appendix analysis turns from the shutdown window to the restoration period. Using a fixed sample of 20,000 Bangladesh-associated accounts selected on pre-shutdown activity, and a text classifier that assigns repositories to NAICS sectors from historical README files, I trace observable push activity by the economic sector associated with the repositories in which it occurs. Restoration of connectivity did not imply an immediate or uniform recovery. Activity in Finance- and Information-related repositories began a two-week run at or above 80 percent of its pre-shutdown level by the third week after the shutdown; Education and Health recovered more gradually; and Accommodation/Food remained well below its baseline throughout the observed window. Activity in pre-existing repositories recovered far less than sector-wide activity, consistent with part of the recovery occurring through new projects. Because this analysis compares each sector with its own pre-shutdown level rather than with an untreated control, it describes heterogeneity in recovery dynamics and does not estimate sector-specific causal effects.

The findings relate to three literatures. First, they add micro-founded evidence to work on the economic costs of Internet shutdowns and, more broadly, of infrastructure outages. Existing studies document broad effects on output, employment, investment, and economic activity \citep*{tagat2024,gunsilius2023}, and \citet{allcott2016} show how firms cope when electricity becomes unavailable. Second, the paper contributes to research on Internet access as a productive input. \citet{hjort2019}, for example, find that the arrival of fast Internet increased employment and firm performance in Africa, and \citet{akerman2015} show that broadband complements skilled work. Third, the paper joins a growing literature that uses public GitHub activity as a high-frequency measure of software work. \citet{kreitmeir2024} exploit Italy's temporary ChatGPT ban and construct daily GitHub outcomes; \citet{hu2023} use open-source contributions to study lockdowns; and \citet{wachs2022} demonstrate the usefulness of GitHub data for measuring the geography of software production.

The interpretation is deliberately narrower than ``total software productivity.'' Public GitHub events omit private repositories, other hosting services, meetings, local coding that is never pushed publicly, and non-coding work. User location is self-reported rather than physically verified. Most importantly, the shutdown was implemented amid severe domestic unrest, so the design cannot separate every contemporaneous Bangladesh-specific shock from the loss of connectivity. The paper therefore estimates the collapse in observable GitHub activity during the shutdown episode, and uses commit timestamps to assess how much of that collapse reflects delayed publication rather than forgone code creation.

The paper proceeds as follows. \Cref{sec:background} sets out the shutdown chronology and explains why GitHub data separate code creation from publication. \Cref{sec:literature} reviews related work, and \cref{sec:data,sec:strategy} describe the sample, outcomes, and estimators. \Cref{sec:results} reports the national, secondary-outcome and international repository-domain synthetic controls alongside the difference-in-differences comparison; discontinuity and descriptive recovery details are in the appendix. \Cref{sec:mechanism} turns to the commit-timestamp evidence, \cref{sec:robustness} examines diagnostics and identification threats, including the concurrent unrest, and \cref{sec:discussion,sec:conclusion} interpret the findings.
